\chapter{Introducción}
\label{cap:intro}

Un ciclo hamiltoniano es un recorrido que pasa por todos los vértices de un grafo una sola vez y
regresa al punto de partida. Saber si un grafo tiene uno parece sencillo, pero es un problema
NP-completo \cite{de1}: no se conoce ningún método que lo resuelva rápido en todos los casos, y
probar todas las rutas posibles deja de ser viable en cuanto el grafo crece un poco (con 100
vértices ya hay más rutas que átomos en el universo). Además, es la base del famoso problema del
agente viajero (TSP) \cite{sa3}.

Por eso en esta práctica usamos \emph{metaheurísticas}: algoritmos que no garantizan encontrar la
respuesta, pero que van mejorando rutas poco a poco hasta acercarse a una solución. Implementamos
seis:

\begin{itemize}
  \item \textbf{Algoritmo genético (GA)}, que usamos como punto de comparación.
  \item \textbf{Recocido simulado (SA)}.
  \item \textbf{Colonia de hormigas (ACO)}.
  \item \textbf{Colonia artificial de abejas (ABC)}.
  \item \textbf{Evolución diferencial (DE)}.
  \item \textbf{Enjambre de partículas (PSO)}.
\end{itemize}

Para probarlos usamos el \textbf{FHCP Challenge Set}, un conjunto de 1\,001 grafos publicado por
el Flinders Hamiltonian Cycle Project \cite{aco8,aco9}. Todos tienen ciclo hamiltoniano, pero
están diseñados a propósito para engañar a las heurísticas sencillas, así que son una buena prueba
para nuestros algoritmos.

Todos los algoritmos comparten el mismo lector de grafos, la misma forma de medir qué tan buena es
una ruta y el mismo verificador final. Así, si un algoritmo funciona mejor que otro, es por la
forma en que busca y no por diferencias en el resto del código.

Hicimos tres experimentos: una comparación principal entre los seis algoritmos (\ejecA{}
corridas), un benchmarking sobre los primeros \grafosB{} grafos del FHCP (\ejecB{} corridas) y una
prueba con grafos grandes, de hasta 9\,528 vértices.

En resumen, los dos algoritmos que mejor funcionaron fueron \alg{ABC} y \alg{ACO} con rotaciones
de Pósa, porque aprovechan la estructura del grafo al buscar. DE y PSO, que trabajan con números
reales en lugar de rutas, no encontraron ningún ciclo. También vimos que los grafos donde todos
los vértices tienen exactamente tres vecinos fueron casi imposibles para todos los algoritmos.

\section*{Cómo está organizado el reporte}
En la sección~\ref{cap:grafo} explicamos qué es un ciclo hamiltoniano y cómo guardamos el grafo.
En la~\ref{cap:repr} vemos cómo representamos las rutas y cómo las calificamos. La~\ref{cap:alg}
describe cada algoritmo. La~\ref{cap:diseno} cuenta cómo hicimos los experimentos y cómo
comprobamos que los resultados fueran correctos. Las secciones~\ref{cap:comp}
y~\ref{cap:bench} muestran los resultados, y la~\ref{cap:concl} cierra con las conclusiones. Al
final, el anexo~\ref{anx:codigo} explica cómo ejecutar el código.
